

\subsection*{Assumption Used}

\begin{assumption}[Documented Scientific Breakthroughs (A5)]\label{A5}
  The following AI-driven scientific results are documented facts.
  For each, A5 explicitly specifies the scientific \emph{field},
  the \emph{type of result}, and the \emph{comparison to the prior
  human record}:
  \begin{enumerate}
    \item \textbf{AlphaFold2} (structural biology, 2021):
      Solved protein-structure prediction to near-experimental accuracy.
      A5 explicitly states this problem \emph{had been open for 50 years}
      despite active human research.  No prior human solution existed.
    \item \textbf{AI drug discovery} (pharmacology, 2023):
      An AI-designed drug candidate reached IND approval in 18 months.
      The historical human-only average is 4--6 years (48--72 months).
      A5 specifies both the AI time and the human-only baseline.
    \item \textbf{GNoME} (materials science, 2023):
      Produced 2.2 million stable crystal structures described as novel
      (not previously in the human record).
      A5 explicitly states this surpasses \emph{the total of all prior
      human experimental discovery combined} by a ratio of $>100\times$.
    \item \textbf{Lean4 AI} (mathematics, 2023--2024):
      Produced novel proof strategies for open combinatorics problems.
      A5 notes such tasks previously required years of specialist effort.
  \end{enumerate}
  \textbf{Domain note (stated in A5):} The four cases belong to four
  different scientific fields as named above: structural biology,
  pharmacology, materials science, and mathematics.  These names appear
  in the documentation of each result and are taken from A5 directly.
\end{assumption}

\subsection*{Definitions}

\begin{definition}[Type-A Beyond-Record Contribution]\label{defA}
  A result $R$ in field $F$ is a \emph{Type-A beyond-record contribution}
  if A5 explicitly states that $R$ solved a problem that had been
  \emph{open} (no prior solution in the human record) in $F$.
\end{definition}

\begin{definition}[Type-B Beyond-Record Contribution]\label{defB}
  A result $R$ in field $F$ is a \emph{Type-B beyond-record contribution}
  if A5 explicitly states that the collection of results produced by $R$
  surpasses \emph{the total of all prior human experimental discovery in
  $F$ combined} by a stated ratio $> 1$.
  (``Experimental discovery'' here follows A5's own phrasing, which
  describes the prior human corpus in materials science as the
  experimental record.  Non-experimental theoretical work in $F$, if
  any, is not separately enumerated in A5 and is therefore not part of
  the comparison basis stated in A5.)
\end{definition}

\begin{definition}[Time-to-Discovery Speedup]\label{defS}
  For a specific discovery task $T$, the \emph{time-to-discovery speedup}
  is the ratio $\tau_{\mathrm{human}}(T) / \tau_{\mathrm{AI}}(T)$, where
  $\tau_{\mathrm{human}}(T)$ is the human-only time to complete $T$ and
  $\tau_{\mathrm{AI}}(T)$ is the AI-assisted time to complete $T$, both
  as stated in A5 for the same task $T$.
\end{definition}

\begin{definition}[Independent Fields]\label{defI}
  Scientific fields are \emph{independent} if each is a distinct named
  discipline and no result in one field constitutes a result in another.
  Structural biology, pharmacology, materials science, and mathematics
  are stipulated as independent: they are standardly classified as
  distinct branches of science, and no protein-structure result is a
  drug-discovery, crystal-structure, or mathematics result.
\end{definition}

\subsection*{Statement to Prove (C3)}

\begin{theorem}[Scientific Force Multiplication]\label{thm:C3}
  Under Assumption~\ref{A5} and
  Definitions~\ref{defA}--\ref{defI}:
  \begin{enumerate}
    \item[(i)] In at least two independent scientific fields, AI has made
      a beyond-record contribution (Type-A or Type-B).
    \item[(ii)] A5 documents at least one time-to-discovery speedup
      $>2\times$, and at least one discovery-scale ratio $>100\times$.
    \item[(iii)] The AI results documented in A5 span at least four
      independent scientific fields (Definition~\ref{defI}).
  \end{enumerate}
\end{theorem}

\begin{proof}

\noindent\textbf{Step 1: Four independent fields — claim (iii).}

A5 names four fields explicitly for the four AI results:
structural biology (AlphaFold2),
pharmacology (AI drug discovery),
materials science (GNoME),
mathematics (Lean4 AI).
By Definition~\ref{defI}, these four fields are independent.
Each case study in A5 constitutes an AI scientific result in its named
field (by A5's own framing).
Therefore, the four AI results documented in A5 span four independent
scientific fields.  Claim (iii) is established.

\medskip
\noindent\textbf{Step 2: Beyond-record contributions in two fields — claim (i).}

\emph{AlphaFold2 — Type-A (Definition~\ref{defA}).}
A5 explicitly states that protein-structure prediction was an
\emph{open problem for 50 years}.  AlphaFold2 solved it.  This is
precisely the condition in Definition~\ref{defA}: A5 states the problem
was open with no prior solution.  AlphaFold2 is a Type-A beyond-record
contribution in structural biology.

\emph{GNoME — Type-B (Definition~\ref{defB}).}
A5 explicitly states that GNoME produced 2.2 million novel crystal
structures and that this \emph{surpasses the total of all prior human
experimental discovery combined} by $>100\times$.  This is precisely
Definition~\ref{defB}: A5 states the collection surpasses the entire
prior human corpus in materials science by a stated ratio ($>100\times$).
GNoME is a Type-B beyond-record contribution in materials science.

Note: Type-A and Type-B are \emph{different} conditions; GNoME is
classified under Type-B, not Type-A (A5 does not say crystal-structure
discovery was an open problem; it says GNoME exceeded the entire prior
corpus in scale).

Two independent fields (structural biology, materials science) exhibit
beyond-record contributions.  Claim (i) is established.

\medskip
\noindent\textbf{Step 3: Quantified acceleration — claim (ii).}

\emph{Time-to-discovery speedup $>2\times$.}
Let task $T_{\mathrm{drug}}$ be the task ``design a drug candidate and
advance it to IND approval.''  A5 specifies both times for this same task:
\[
  \tau_{\mathrm{human}}(T_{\mathrm{drug}}) \;\geq\; 4\text{ years}
    = 48\text{ months}
  \quad\text{(A5: historical human-only average for } T_{\mathrm{drug}}\text{)},
\]
\[
  \tau_{\mathrm{AI}}(T_{\mathrm{drug}}) = 18\text{ months}
  \quad\text{(A5: Insilico Medicine result for } T_{\mathrm{drug}}\text{)}.
\]
By Definition~\ref{defS}, using the conservative baseline:
\[
  \text{Speedup}(T_{\mathrm{drug}})
  = \frac{\tau_{\mathrm{human}}}{\tau_{\mathrm{AI}}}
  \geq \frac{48}{18} = \frac{8}{3} \approx 2.67 > 2.
\]

\emph{Discovery-scale ratio $>100\times$.}
A5 explicitly states the GNoME ratio is $>100\times$.
We adopt this directly from A5 without further arithmetic.

Both quantified thresholds in claim (ii) are established.

\medskip
\noindent\textbf{Conclusion.}
Claims (i), (ii), and (iii) are derived from A5 and
Definitions~\ref{defA}--\ref{defI} only.
AI is a cross-domain scientific force multiplier with beyond-record
contributions in $\geq 2$ independent fields, a confirmed $>2\times$
time speedup, a confirmed $>100\times$ scale ratio, spanning 4
independent scientific disciplines.
\end{proof}

